%%%%%%%%%%%%%%%%%%%%%%%%%%%%%%%%%%%%%%%%%%%%%%%%%%%%%%%%%%%%%%%%%%%%%%%%%%%%%%%%%%%%%%%%%%%%%%%%%
% Abstract   (the \chapter* command is in main.tex)
%%%%%%%%%%%%%%%%%%%%%%%%%%%%%%%%%%%%%%%%%%%%%%%%%%%%%%%%%%%%%%%%%%%%%%%%%%%%%%%%%%%%%%%%%%%%%%%%%

A delivery drone must decide when to abandon its route and fly home, yet it cannot measure its battery's remaining charge directly: its current sensor is noisy and biased, and its voltage sags under load. This report develops an Extended Kalman Filter (EKF) that estimates the state of charge (SoC) and its uncertainty from these two signals, and a risk-aware return-to-home rule that turns back when the probability of landing with a 10\% reserve falls below 95\%. Both are tested in a MATLAB simulator whose battery deliberately differs from the filter's equivalent-circuit model. A design-stage investigation showed that resistance mismatch dominated the SoC error, so the filter learns the cell resistance during flight; this reduced the SoC RMSE on a nominal flight to 0.97\%, against 4.36\% for charge counting. A five-state filter that also estimates sensor bias and capacity kept honest uncertainty bands in every nominal, biased and aged-battery scenario and was chosen for decisions. In 500 random missions, the risk-aware rule had no battery failures and broke the reserve in 2.8\% of missions, against 38.0\% for the same filter without its uncertainty and 30.4\% battery failures for charge counting. Its breach rate stayed within the promised level at every threshold tested, showing that its probabilities were calibrated.
